\scanpage{137}
\mspage{82}
% archive sheet number: 136

pouvoir s'en débrouiller. Mais la formule d'adjonction me semble rester inattendue et profonde, même en admettant les résultats de Quillen \inserted{et} Thomason. Elle paraît liée \inserted{ici} à \deleted{une} propriété assez particulière du foncteur $C$, savoir le lemme p.~68

% slanted
\margin{Voir \unsure{après coup|ap. ci-après} le th. sch. \unsure{2, -- 2°}, p 108.113. Ce n'est finalement pas si prof\unsure{ond} qu'il m'en semblait, mais quasiment sorital!}

\[
\begin{aligned}
,\,i_{C(X)}&=C(i_X):C(X)\unsure{\mathrel{\relbar\joinrel\relbar\joinrel\vert}|\longrightarrow}\unsure{C^2(C(X))|C(C(X))}\\
&\overset{\uparrow}{\text{dans }}\underbrace{(\mathrm{Cat}_{/S})(W_S^d)^{-1}}_{\mathrm{HOT}_W(S)}
 \text{ et dans }\underbrace{(\mathrm{Cat}_{/S})(W_S^g)^{-1}}_{\mathrm{HOT}_W^{\unsure{\circ|o}}(S)},
\end{aligned}
\]

ce qui n'a pas l'air purement formel. \deleted{\unsure{Il le semble}} Il serait intéressant d'examiner dans quelle mesure \deleted{elle} \deleted{\unsure{pourrait}} cette formule garde un sens \deleted{pour} dans une ``catégorie de modèles'' $(\mathcal M,W)$ plus ou moins quelconque, \deleted{\unsure{du moins en}} ou provenant d'un triple de Quillen $(W,\mathrm{Cof},\mathrm{Fib})$ dans $\mathcal M$. Il est clair que la factorisation $X\overset{f}{\unsure{\mathrel{\relbar\joinrel\relbar\joinrel\vert}|\longrightarrow}}S$ en $X\overset{i_X}{\unsure{-|\longrightarrow}}C(X)\overset{\bar f}{\unsure{-|\longrightarrow}}S$ ($i_X\in W$, $\bar f\in\mathrm{Parf}$, jouant le rôle des ``fibrations'' dans le contexte Quillen, a un sens intrinsèque indépendt du choix d'une construction particulière, chose qu'on peut préciser de façon parfaite. C'est la \deleted{\unsure{réd.}} description \unsure{des localiseurs|du localiseur} $W_S^d,W_S^g$ dans $\mathcal M_{/S}$ (pour $S\in\mathrm{Ob}\,\mathcal M$) qui n'est pas comprise de façon générale, en dehors du cas $\mathcal M=\mathrm{Cat}$. En d'autres termes, je ne comprends pas comment construire des catégories $\mathrm{\unsure{HOt|HOT}}_{\mathcal M,W}(S)$, $\mathrm{HOT}_{\mathcal M,W}^{\circ}(S)$, $=$ si $(\mathcal M,W)$ provient d'un triple de Quillen, j'\deleted{\unsure{ai|avais}} \unsure{en|n'}\inserted{vois} seulement la description d'une catégorie $\mathrm{HOT}^{\unsure{\mathrm{lc}|\mathrm{le}}}_{\mathcal M,W}(S)$ \inserted{$=(\mathcal M_{/S})(W_S^{\mathrm{grossier}})^{-1}$}, et un foncteur \inserted{de localisation}

\[
\mathcal M_{/S}\longrightarrow\mathrm{HOT}_{\mathcal M,W}^{\mathrm{lc}}(S).
\]

Il faut arriver à décrire des \inserted{$\mathcal M_{/S}$ \deleted{\unsure{des}}} $W$-équivalences (``\uline{relatives à $S$}'') plus fines que la grossière, \deleted{\unsure{induite}} \inserted{donnée} par $W$ sur la source $X$ sans plus. Dans $\mathrm{Cat}$ j'y arrive en utilisant \inserted{soit} la famille des isom. locaux (pour décrire $W_S^{\deleted{\unsure{g}}d}$), soit celle des isom colocaux (pour décrire $W_S^g$) -- en tous les cas il y a une notion de ``localisation'' (soit sur $S$, soit sur $S^\circ$\ldots), ou
